%!TEX root = ../main.tex

\section{Полная система уравнений модели}

Перечислим уравнения модели. Для начала, это уравнения Максвелла в электростатическом приближении \eqref{eq:electrostatic}.

Плотность свободной энергии среды после конкретизации общего вида \eqref{eq:free_energy_density} определяется выражением
\[
    \psi = -\half \epsilon[\phi] \vE \cdot \vE + \theta \left[ c(\phi) - c(\phi) \ln \cfrac{\theta}{\theta_0} - \eta_0(\phi) \right] + \cfrac{\Gamma}{l^2}[1 - f(\phi)] + \cfrac{\kappa}{2} \rho^2 + \cfrac{\Gamma}{4} \nabla \phi \cdot \nabla \phi.
\]

Уравнение \eqref{eq:temperature} изменения температуры принимает форму
\[
    c(\phi) \dt{\theta} = - \Div(-k \nabla \theta) + \cfrac{1}{m} \left( \partt{\phi} \right)^2 + \sigma[\phi] (\vE - \kappa \nabla \rho) \cdot (\vE - \kappa \nabla \rho) - \theta \partt{\phi} \left[ c'_\phi \ln \cfrac{\theta}{\theta_0} + (\eta_0)'_\phi \right].
\]

Уравнение \eqref{eq:phi} динамики фазового поля -- форму
\[
    \cfrac{1}{m} \partt{\phi} = \half \epsilon'_\phi \vE \cdot \vE + \theta \left[ -c'_\phi + c'_\phi \ln \cfrac{\theta}{\theta_0} + (\eta_0)'_\phi \right] + \cfrac{\Gamma}{l^2} f'_\phi + \cfrac{\Gamma}{2} \Delta \phi.
\]

Диэлектрическая проницаемость $\epsilon(\vx, \phi) = \epsilon[\phi]$ и электропроводность $\sigma(\vx, \phi) = \sigma[\phi]$ среды заданы формулами \eqref{eq:electrical_properties}. $m$, $k$, $\kappa$, $\Gamma$, $l$, $\theta_0$ -- константы; $c(\phi)$ и $\eta_0(\phi)$, вообще говоря, произвольные положительные функции, считающиеся известными.

С учетом произведенных уточнений, определяющие соотношения модели имеют вид
{\allowdisplaybreaks
\begin{align*}
    \eta &= c(\phi) \ln \cfrac{\theta}{\theta_0} + \eta_0(\phi) \\
    \vD &= \epsilon[\phi] \vE, \\
    \vxi &= \cfrac{\Gamma}{2} \nabla \phi, \\
    \pi &= \half \epsilon'_\phi \vE \cdot \vE + \theta \left[ -c'_\phi + c'_\phi \ln \cfrac{\theta}{\theta_0} + (\eta_0)'_\phi \right] + \cfrac{\Gamma}{l^2} f'_\phi - \cfrac{1}{m} \partt{\phi}, \\
    \vq &= -k \nabla \theta, \\
    \vj &= \sigma[\phi] (\vE - \kappa \nabla \rho).
\end{align*}
}

Из соотношения \eqref{eq:free_energy_density} выразим плотность внутренней энергии среды. Учитывая, что из формул \eqref{eq:entropy_thermal_free_energy} $\psi^T + \theta \eta = \theta c(\phi)$, получим
\[
    e = \half \epsilon[\phi] \vE \cdot \vE + \theta c(\phi) + \cfrac{\Gamma}{l^2}[1 - f(\phi)] + \cfrac{\kappa}{2} \rho^2 + \cfrac{\Gamma}{4} \nabla \phi \cdot \nabla \phi.
\]